% संपूर्ण मराठी ग्रंथातील प्रकरण 55: अंतःप्रज्ञावादी टॅब्लो
% हा संपादनयोग्य स्रोतखंड आहे. संकलनासाठी संपूर्ण स्रोतसंग्रहातील
% अक्षररूपे, पूर्वभाग, चिन्हव्याख्या आणि आकृती-स्रोत आवश्यक आहेत.
% मूळ: Open Logic Project; मजकूर आणि रूपांतर CC BY 4.0.
% यंत्रानुवाद, दुरुस्ती आणि तपासणी: OpenAI Codex —
% GPT-5.6 Sol आणि GPT-6 Sol, दोन्ही Ultra effort.
% दुरुस्ती-पुनरावलोकन: GPT-6.1 Sol, Ultra effort.
\chapter{अंतःप्रज्ञावादी टॅब्लो}\label{int:tab:chap}
\begin{editorial}
  अंतःप्रज्ञावादी तर्कशास्त्रातील पूर्वचिन्हयुक्त
  टॅब्लोंविषयी हे मसुदा प्रकरण आहे.
  अधिक उदाहरणे, संपूर्णतेच्या सिद्धता,
  आणि बंद टॅब्लो शोधण्याचा प्रयत्न
  अयशस्वी झाल्यावर प्रतिप्रतिमाने कशी
  मिळवता येतात याची चर्चा हवी आहे.
\end{editorial}









\section{प्रस्तावना}\label{int:tab:int:sec}

अभिजात तर्कशास्त्रातील टॅब्लो म्हणजे चिन्हांकित सूत्रांचे (खाली शाखा फुटणारे) काही वृक्ष.
चिन्हांकित सूत्र म्हणजे सत्यतामूल्याचे चिन्ह
($\True$ किंवा $\False$) आणि वाक्य
यांची जोडी:
\[
\sFmla{\True}{\varphi} \text{किंवा} \sFmla{\False}{\varphi}.
\]
टॅब्लोची सुरुवात काही
\emph{गृहीतकां}पासून होते. पुढील प्रत्येक
चिन्हांकित सूत्र हे अनुमानाचा एखादा
नियम लावून मिळते. काही अनुमान-नियम
वृक्षाच्या टोकावर एक किंवा अधिक
चिन्हांकित सूत्रे जोडतात; इतर
नियम दोन नवी टोके जोडून दोन शाखा
तयार करतात. नियम लावल्यावर मिळणाऱ्या
चिन्हांकित सूत्रे मधील
सूत्र, नियम ज्या
चिन्हांकित सूत्राला लावला
त्यातील सूत्रापेक्षा कमी गुंतागुंतीचे असते.
एखाद्या शाखेत
$\sFmla{\True}{\varphi}$ आणि
$\sFmla{\False}{\varphi}$ ही दोन्ही
असतील, तर ती शाखा \emph{बंद}
आहे असे म्हणतो. टॅब्लो
मधील प्रत्येक शाखा बंद असेल,
तर संपूर्ण टॅब्लो बंद असतो.
बंद टॅब्लो ही अशी
निष्पत्ती असते की तिच्यावरून
त्या टॅब्लोची सुरुवात ज्या
चिन्हांकित सूत्रांनी केली तो
संच अपूर्ततायोग्य आहे हे दिसते.
यावरून $\Proves$ संबंधाची व्याख्या
करता येते: $\Gamma \Proves \varphi$
तेव्हा आणि केवळ तेव्हाच
$\Gamma_0 = \{\psi_1,
\dots, \psi_n\} \subseteq \Gamma$
असा एखादा सांत संच असतो की
पुढील गृहीतकांसाठी बंद
टॅब्लो मिळतो:
\[
\{\sFmla{\False}{\varphi}, \sFmla{\True}{\psi_1}, \dots, \sFmla{\True}{\psi_n}\}.
\]

अंतःप्रज्ञावादी तर्कशास्त्रासाठी
चिन्हांकित सूत्र ही संकल्पना
विस्तारित करावी लागते आणि
तार्किक कारकांचे नियम बदलावे लागतात.
चिन्हाबरोबर ($\True$ किंवा
$\False$) अंतःप्रज्ञावादी
टॅब्लो मधील सूत्रे ना
\emph{पूर्वचिन्हे}~$\sigma$ ही असतात.
ही पूर्वचिन्हे धन पूर्णांकांच्या
रिकाम्या नसलेल्या क्रमिका असतात;
म्हणजे $\sigma \in
(\PosInt)^* \setminus \{\emptyseq\}$.
अशी पूर्वचिन्हे लिहिताना
भोवतालचे $\tuple{\ }$ वगळतो
आणि वेगवेगळ्या घटक मध्ये
$,$ ऐवजी $.$ वापरतो.
$\sigma$ हे पूर्वचिन्ह असेल,
तर $\sigma.n$ म्हणजे
$\sigma \concat \tuple{n}$;
उदाहरणार्थ $\sigma = 1.2.1$
असेल, तर $\sigma.3$ म्हणजे
$1.2.1.3$. म्हणून उदाहरणार्थ,
\[
\sFmla{\True}{\varphi \lif (\psi \lif \chi)}[1.2]
\]
हे \emph{पूर्वचिन्हयुक्त
चिन्हांकित सूत्र} (किंवा संक्षेपाने
\emph{पूर्वचिन्हयुक्त सूत्र})
आहे.

सहज अर्थाने, टॅब्लोच्या
एखाद्या शाखेवरील सूत्रे
पूर्ण करू शकेल अशा प्रतिमानातील
जगाचे नाव पूर्वचिन्ह देते.
आणि $\sigma$ हे एखाद्या जगाचे
नाव असेल, तर $\sigma.n$ हे
म्हणजे~$\sigma$ ने दर्शवलेल्या
जगापासून प्राप्य असलेल्या
जगाचे नाव देते.

अंतःप्रज्ञावादी प्रतिमानांमध्ये
प्राप्यता संबंध परावर्ती आणि
संक्रमक असतो. पूर्वचिन्हांच्या
भाषेत याचा अर्थ असा की
$\sigma$ हा स्वतः~$\sigma$ पासून प्राप्य
असतो; आणि~$\sigma$ चा कोणताही
विस्तार, म्हणजे
$\sigma.n_1.\cdots.n_k$
या रूपातील पूर्वचिन्हसुद्धा,
त्यापासून प्राप्य असते.
$\sigma$ आणि त्याचा कोणताही
विस्तार दर्शवण्यासाठी
$\sigma.*$ ही संज्ञा वापरू.
म्हणजे $\sigma.*$ ही पूर्वचिन्हे
नेमकी~$\sigma$ पासून प्राप्य
असलेली सर्व पूर्वचिन्हे आहेत.






\begin{quote}\small\textbf{स्रोतनोंद.} SOL6-A206: उलट चिन्हांची जोडी आणि संपूर्ण शाखेची बंदता यांचे scopes वेगळे केले. नवे साक्षी-पूर्वचिन्ह किंवा त्याचा विस्तार शाखेवर आधी नसावा अशी अट स्पष्ट केली, जेणेकरून आधीच्या जगांचे अर्थनिर्धारण जपता येते.\end{quote}






\section{अंतःप्रज्ञावादी तर्कशास्त्रासाठीचे नियम}\label{int:tab:rul:sec}

$\land$ आणि $\lor$ या कारकांचे नियम नेहमीच्या विधानीय
चिन्हांकित टॅब्लो मधील नियमांसारखेच आहेत;
फक्त त्यांना पूर्वचिन्हे जोडली जातात. प्रत्येक वेळी,
$\sFmla{S}{\varphi}[\sigma]$ या चिन्हांकित सूत्र
ला नियम लावल्याने मिळणाऱ्या नव्या सूत्रे
वरही~$\sigma$ हेच पूर्वचिन्ह असते. हे सहज लक्षात
येते: उदाहरणार्थ, $\varphi \land \psi$ हे~$\sigma$ ने
नाव दिलेल्या जगात सत्य असेल, तर $\varphi$ आणि $\psi$
ही दोन्ही~$\sigma$ येथे सत्य असतात (आणि त्या
जगापासून प्राप्य असलेल्या प्रत्येक जगातही सत्य
असतात). $\land$ आणि $\lor$ यांचे नियम
\ref{int:tab:rul:tab:prop-rules} मध्ये दिले आहेत.

\begin{table}
  \[\def\arraystretch{3}\begin{array}{|c|c|}
    \hline
    \AxiomC{\sFmla{\True}{\varphi \land \psi}[\sigma]}
    \RightLabel{\TRule{\True}{\land}}
    \UnaryInfC{\sFmla{\True}{\varphi}[\sigma]}
    \noLine
    \UnaryInfC{\sFmla{\True}{\psi}[\sigma]}
    \DisplayProof
    &
    \AxiomC{\sFmla{\False}{\varphi \land \psi}[\sigma]}
    \RightLabel{\TRule{\False}{\land}}
    \UnaryInfC{$\sFmla{\False}{\varphi}[\sigma] \quad \mid \quad
      \sFmla{\False}{\psi}[\sigma]$}
    \DisplayProof
    \\[2ex]
    \hline
    \AxiomC{\sFmla{\True}{\varphi \lor \psi}[\sigma]}
    \RightLabel{\TRule{\True}{\lor}}
    \UnaryInfC{$\sFmla{\True}{\varphi}[\sigma] \quad \mid \quad
      \sFmla{\True}{\psi}[\sigma]$}
    \DisplayProof
    &
    \AxiomC{\sFmla{\False}{\varphi \lor \psi}[\sigma]}
    \RightLabel{\TRule{\False}{\lor}}
    \UnaryInfC{\sFmla{\False}{\varphi}[\sigma]}
    \noLine
    \UnaryInfC{\sFmla{\False}{\psi}[\sigma]}
    \DisplayProof
    \\[2ex]
    \hline
  \end{array}\]
  \caption{$\land$ आणि $\lor$ साठीचे पूर्वचिन्हयुक्त
    टॅब्लो नियम}
  \label{int:tab:rul:tab:prop-rules}
\end{table}

शाखा बंद होण्याची अट नेहमीच्या टॅब्लो
प्रमाणेच आहे; पण येथे सूत्रे बरोबर त्यांची
पूर्वचिन्हेही जुळली पाहिजेत. खरे तर ही अट
थोडी शिथिल करता येते: अंतःप्रज्ञावादी प्रतिमानात
सूत्रे एकदा सत्य झाली की पुढेही सत्य राहतात.
म्हणून $\varphi$ हे~$\sigma$ येथे सत्य असताना
ते~$\sigma.{*}$ या प्राप्य पूर्वचिन्हावर असत्य
असू शकत नाही. म्हणून एखाद्या शाखेत पुढील
दोन्ही असतील, तर ती बंद असते:
\[
\sFmla{\True}{\varphi}[\sigma] \quad\text{आणि}\quad \sFmla{\False}{\varphi}[\sigma.{*}]
\]
येथे $\sigma$ हे एखादे पूर्वचिन्ह आणि $\varphi$ हे
सूत्र आहे. उलट चिन्हे असतील, म्हणजे
शाखेत पुढील दोन्ही असतील,
\[
\sFmla{\False}{\varphi}[\sigma] \quad\text{आणि}\quad \sFmla{\True}{\varphi}[\sigma.{*}]
\]
तर सदर सूत्र-जुळणीची अट या उलट चिन्हांच्या जोडीला
तेव्हा आणि केवळ तेव्हाच लागू होते, जेव्हा $*$ ही
रिकामी क्रमिका असते. शाखा इतर कोणत्या बंदतेच्या
अटीमुळे बंद होऊ शकते, हे याने नाकारलेले नाही.

तसेच, एखाद्या शाखेत~$\sFmla{\True}{\bot}[\sigma]$
असेल, तरी ती बंद असते.

गृहीतके मांडण्याचे नियमही नेहमीच्या टॅब्लो
प्रमाणेच आहेत; मात्र गृहीतकांसाठी आपण नेहमी
$1$ हे पूर्वचिन्ह वापरतो. (सर्व गृहीतकांसाठी
एकच पूर्वचिन्ह वापरले, तर नेमके कोणते
पूर्वचिन्ह वापरले याने फरक पडत नाही.)
म्हणून, उदाहरणार्थ,
\[
\psi_1, \dots, \psi_n \Proves \varphi
\]
असे तेव्हा आणि केवळ तेव्हाच म्हणतो जेव्हा पुढील
गृहीतकांसाठी बंद टॅब्लो असतो:
\[
\sFmla{\True}{\psi_1}[1], \dots, \sFmla{\True}{\psi_n}[1],
\sFmla{\False}{\varphi}[1].
\]

शर्त~$\lif$ साठीचे नियम अभिजात आणि मोडल
प्रकरणांपेक्षा वेगळे आहेत. $\TRule{\True}{\lif}$
हा नियम $\sFmla{\True}{\varphi \lif \psi}[\sigma]$
असलेल्या शाखेतून दोन वेगळ्या शाखा काढतो:
एकीवर $\sFmla{\False}{\varphi}[\sigma.{*}]$ आणि
दुसरीवर $\sFmla{\True}{\psi}[\sigma.{*}]$ जोडतो.
हा नियम लावताना $\sigma.{*}$ हे पूर्वचिन्ह
संबंधित शाखेत \emph{आधीपासूनच} आलेले
असले पाहिजे. अशा पूर्वचिन्हाला त्या शाखेवर
‘वापरलेले’ म्हणू. ($\sigma.{*}$ मध्ये
$\sigma$ स्वतःही येत असल्याने, वेगवेगळ्या
शाखांवर $\sFmla{\False}{\varphi}[\sigma]$ आणि
$\sFmla{\True}{\psi}[\sigma]$ जोडून हा नियम
नेहमी लावता येतो.)

$\TRule{\False}{\lif}$ हा नियम
$\sFmla{\False}{\varphi \lif \psi}[\sigma]$
असलेल्या शाखेवर
$\sFmla{\True}{\varphi}[\sigma.n]$ आणि
$\sFmla{\False}{\psi}[\sigma.n]$ ही दोन्ही
सूत्रे जोडतो; येथे $\sigma.n$ हे त्या
शाखेसाठी नवे पूर्वचिन्ह असते: ते किंवा त्याचा कोणताही
विस्तार त्या शाखेवर आधी आलेला नसतो. पुढील नकरणाच्या
नियमांतही ‘नवे’ याच अर्थाने वापरतो.

$\lnot$ साठीचे नियम याच धर्तीवर व्याख्यायित
केले आहेत ($\lnot \varphi$ म्हणजे
$\varphi \lif \lfalse$ ही व्याख्या वापरून).

हे नियम \ref{int:tab:rul:tab:rules-lif-lnot} मध्ये दिले आहेत.

\begin{table}
  \[\def\arraystretch{3}\begin{array}{|c|c|}
    \hline
    \AxiomC{\sFmla{\True}{\lnot \varphi}[\sigma]}
    \RightLabel{\TRule{\True}{\lnot}}
    \UnaryInfC{$\sFmla{\False}{\varphi}[\sigma.{*}]$}
    \DisplayProof
    &
    \AxiomC{\sFmla{\False}{\lnot \varphi}[\sigma]}
    \RightLabel{\TRule{\False}{\lnot}}
    \UnaryInfC{\sFmla{\True}{\varphi}[\sigma.n]}
    \DisplayProof
    \\[1ex]
    \text{$\sigma.{*}$ वापरलेले आहे} & \text{$\sigma.n$ नवे आहे}\\
    \hline
    \AxiomC{\sFmla{\True}{\varphi \lif \psi}[\sigma]}
    \RightLabel{\TRule{\True}{\lif}}
    \UnaryInfC{$\sFmla{\False}{\varphi}[\sigma.{*}] \quad \mid \quad
      \sFmla{\True}{\psi}[\sigma.{*}]$}
    \DisplayProof
    &
    \AxiomC{\sFmla{\False}{\varphi \lif \psi}[\sigma]}
    \RightLabel{\TRule{\False}{\lif}}
    \UnaryInfC{\sFmla{\True}{\varphi}[\sigma.n]}
    \noLine
    \UnaryInfC{\sFmla{\False}{\psi}[\sigma.n]}
    \DisplayProof
    \\[1ex]
    \text{$\sigma.{*}$ वापरलेले आहे} & \text{$\sigma.n$ नवे आहे}\\
    \hline
  \end{array}\]
  \caption{$\lnot$ आणि $\lif$ साठीचे
    पूर्वचिन्हयुक्त टॅब्लो नियम}
  \label{int:tab:rul:tab:rules-lif-lnot}
\end{table}












\section{अंतःप्रज्ञावादी तर्कशास्त्रासाठीचे टॅब्लो}\label{int:tab:prf:sec}

\begin{ex}
  $(\varphi \land \psi) \lif \chi \Proves
  \varphi \lif (\psi \lif \chi)$ हे दाखवणारा
  बंद टॅब्लो पुढे दिला आहे.
  \begin{oltableau}
    [\pFmla{\True}{(\formula{A} \land \formula{B}) \lif \formula{C}}{1},
      just =\TAss
      [\pFmla{\False}{\formula{A} \lif (\formula{B} \lif \formula{C})}{1},
        just =\TAss
        [\pFmla{\True}{\formula{A}}{1.1},
          just = {\TRule{\False}{\lif}[2]}
          [\pFmla{\False}{\formula{B} \lif \formula{C}}{1.1},
            just = {\TRule{\False}{\lif}[2]}
            [\pFmla{\True}{\formula{B}}{1.1.1},
              just = {\TRule{\False}{\lif}[4]}
              [\pFmla{\False}{\formula{C}}{1.1.1},
                just = {\TRule{\False}{\lif}[4]}
                [\pFmla{\False}{\formula{A} \land \formula{B}}{1.1.1},
                  just = {\TRule{\True}{\lif}[1]}
                  [\pFmla{\False}{\formula{A}}{1.1.1},
                    just= {\TRule{\False}{\land}[7]}, close]
                  [\pFmla{\False}{\formula{B}}{1.1.1},
                    just= {\TRule{\False}{\land}[7]}, close]]
                [\pFmla{\True}{\formula{C}}{1.1.1},
                  just = {\TRule{\True}{\lif}[1]}, close]]
            ]
          ]
        ]
      ]
    ]
  \end{oltableau}
\end{ex}



\begin{prob}
  पुढील निष्कर्ष दाखवण्यासाठी बंद अंतःप्रज्ञावादी
  टॅब्लो शोधा:
  \begin{enumerate}
    \item $\Proves \varphi \lif (\psi \lif \varphi)$
    \item $\Proves \lnot(\varphi \land \lnot \varphi)$
    \item $\varphi \lif (\psi \lif \chi) \Proves (\varphi \land \psi) \lif \chi$
    \item $\lnot \varphi \lor \lnot \psi \Proves \lnot(\varphi \land \psi)$
  \end{enumerate}
\end{prob}






\begin{quote}\small\textbf{स्रोतनोंद.} SOL6-A207: मूळ adjacent-only अर्थनिर्धारणातील gap दुरुस्त करून सर्व वापरलेल्या पूर्वचिन्ह-विस्तारांसाठी प्राप्यता आवश्यक केली. नवा witness prefix घेताना सर्व जुने ancestors आणि जगांचे अर्थनिर्धारण जपले जाते हे सिद्ध केले; OLINC-233 ही ऐतिहासिक स्रोतदोषाची नोंद जपली.\end{quote}






\section{अंतःप्रज्ञावादी टॅब्लोची निर्दोषता}\label{int:tab:sou:sec}

\begin{explain}
  अंतःप्रज्ञावादी टॅब्लो निर्दोष आहेत हे
  दाखवण्यासाठी पुढील गृहीतकांना
  \[
  \sFmla{\True}{\psi_1}[1], \dots, \sFmla{\True}{\psi_n}[1], \sFmla{\False}{\varphi}[1]
  \]
  बंद टॅब्लो असेल, तर $\psi_1, \dots, \psi_n \Entails \varphi$
  हे दाखवावे लागेल. व्यत्यस्त रूप सिद्ध करणे सोपे आहे:
  एखाद्या $\mModel{M}$ प्रतिमानात आणि जगात~$w$,
  सर्व $i=1$, \dots,~$n$ साठी
  $\mSat{M}{\psi_i}[w]$ असेल, पण
  $\mSat/{M}{\varphi}[w]$ असेल, तर कोणताही
  टॅब्लो बंद होऊ शकत नाही. असे
  प्रतिप्रतिमान टॅब्लोची सुरुवातीची
  गृहीतके पूर्ततायोग्य असल्याचे दाखवते.
  पुराव्याची युक्ती अशी: टॅब्लो
  च्या शाखेवरील सर्व पूर्वचिन्हयुक्त
  सूत्रे पूर्ततायोग्य असतील, तर
  कोणताही नियम लावल्यानंतर मिळणाऱ्या
  विस्तारित शाखांपैकी किमान एक शाखा
  पूर्ततायोग्य असते. बंद शाखा अपूर्ततायोग्य
  असल्याने, पूर्ततायोग्य पूर्वचिन्हयुक्त
  सूत्रांच्या संचावरील कोणत्याही
  टॅब्लोमध्ये किमान एक खुली शाखा
  असलीच पाहिजे.

  ही युक्ती अंतःप्रज्ञावादी प्रकरणात
  वापरण्यासाठी ‘पूर्ततायोग्य’ या संकल्पनेची
  व्याख्या प्राप्यता-संबंध आणि पूर्वचिन्हे
  यांना सामावून घेईल अशी वाढवावी लागते.
  त्यानंतर पुरावा सरळ आहे.
\end{explain}

\begin{defn}
  $P$ हा पूर्वचिन्हांचा एखादा संच असू द्या,
  म्हणजे $P \subseteq (\PosInt)^*
  \setminus \{\emptyseq\}$; तसेच
  $\mModel{M}$ हे प्रतिमान असू द्या.
  $f\colon P\to W$ हे फलन $\mModel{M}$ मधील
  $P$ चे \emph{अर्थनिर्धारण} असते तेव्हा आणि केवळ तेव्हाच,
  जेव्हा प्रत्येक $\sigma,\tau\in P$ साठी $\sigma$ हा
  $\tau$ चा आरंभीचा खंड असेल, तर $Rf(\sigma)f(\tau)$.
  ही अट सर्व विस्तारांसाठी आहे; मधली पूर्वचिन्हे
  $P$ मध्ये असणे आवश्यक नाही.

  पूर्वचिन्हे~$P$ यांच्या अर्थनिर्धारणाच्या
  संदर्भात पुढील व्याख्या करता येतात:
  \begin{enumerate}
  \item $\mModel{M}$ हे $\sFmla{\True}{\varphi}[\sigma]$
    पूर्ण करते तेव्हा आणि केवळ तेव्हाच
    $\mSat{M}{\varphi}[f(\sigma)]$.
  \item $\mModel{M}$ हे $\sFmla{\False}{\varphi}[\sigma]$
    पूर्ण करते तेव्हा आणि केवळ तेव्हाच
    $\mSat/{M}{\varphi}[f(\sigma)]$.
  \end{enumerate}
\end{defn}

\begin{quote}\small\textbf{स्रोतनोंद.} OLINC-233: मूळ adjacent-prefix व्याख्येमुळे मधली पूर्वचिन्हे नसताना दीर्घ विस्तारापर्यंत प्राप्यता मिळत नव्हती. सध्याच्या मराठी व्याख्येत सर्व पूर्वचिन्ह-विस्तारांसाठी प्राप्यता थेट आवश्यक केली आहे; खालील बंदता-सिद्धतेला आता ती अट उपलब्ध आहे.\end{quote}
या व्याख्येनुसार $\sigma$ आणि $\sigma.{*}$ ही दोन्ही
पूर्वचिन्हे वापरलेली असतील, तर $Rf(\sigma)f(\sigma.{*})$.
रिकाम्या विस्तारासाठी हे $R$ च्या परावर्तितेनेही मिळते.

\begin{defn}
  $\Gamma$ हा पूर्वचिन्हयुक्त सूत्रे
  चा संच असू द्या आणि त्यात आलेल्या
  पूर्वचिन्हांचा संच $P(\Gamma)$ असू द्या.
  $f$ ही $P(\Gamma)$ ची
  $\mModel{M}$ मधील अर्थनिर्धारण
  असेल, तर $\mModel{M}$ हे
  $f$ च्या संदर्भात $\Gamma$
  पूर्ण करते, म्हणजे
  $\mSat{M}{\Gamma}[f]$, असे
  तेव्हा म्हणतो जेव्हा
  $\mModel{M}$ हे $\Gamma$ मधील
  प्रत्येक पूर्वचिन्हयुक्त
  सूत्र हे~$f$ च्या संदर्भात
  पूर्ण करते. $\mModel{M}$ असे
  प्रतिमान आणि $P(\Gamma)$ ची
  $f$ असे अर्थनिर्धारण अस्तित्वात
  असतील की $\mSat{M}{\Gamma}[f]$,
  तेव्हा आणि केवळ तेव्हाच $\Gamma$ \emph{पूर्ततायोग्य}
  आहे.
\end{defn}

\begin{prop}
  $\Gamma$ मध्ये एखाद्या
  सूत्र~$\varphi$ आणि पूर्वचिन्ह~$\sigma$
  साठी $\sFmla{\True}{\varphi}[\sigma]$ आणि
  $\sFmla{\False}{\varphi}[\sigma.{*}]$
  ही दोन्ही असतील, किंवा
  $\sFmla{\True}{\lfalse}[\sigma]$
  असेल, तर $\Gamma$ अपूर्ततायोग्य आहे.
\end{prop}

\begin{proof}
  नेहमीच $\mSat/{M}{\lfalse}[f(\sigma)]$
  असल्याने, $\sFmla{\True}{\lfalse}$
  असलेला $\Gamma$ अपूर्ततायोग्य असतो.

  आता दोन्ही चिन्हांकित सूत्रे पूर्ण
  करणारे $\mModel{M}$ प्रतिमान आणि
  $P(\Gamma)$ चे $f$ अर्थनिर्धारण
  असू शकत नाहीत. कारण
  $\mSat{M}{\varphi}[f(\sigma)]$
  असेल, तर
  \ref{int:sem:rel:prop:true-monotonic}
  आणि $Rf(\sigma)f(\sigma.{*})$
  यांवरून
  $\mSat{M}{\varphi}[f(\sigma.{*})]$
  मिळते. म्हणून
  $\mSat{M}{\varphi}[f(\sigma)]$ आणि
  $\mSat/{M}{\varphi}[f(\sigma.{*})]$
  ही दोन्ही एकत्र असू शकत नाहीत.
\end{proof}

\begin{thm}[निर्दोषता]
  \label{int:tab:sou:thm:tableau-soundness}
  $\Gamma$ ला बंद टॅब्लो
  असेल, तर $\Gamma$ अपूर्ततायोग्य
  आहे.
\end{thm}

\begin{proof}
टॅब्लोच्या एखाद्या शाखेवरील
चिन्हांकित सूत्रांचा संच
पूर्ततायोग्य असेल, तर ती शाखा
पूर्ततायोग्य म्हणू. किमान एक
पूर्ततायोग्य शाखा असलेला
टॅब्लो पूर्ततायोग्य म्हणू.

आपण पुढील विधान दाखवू:
पूर्ततायोग्य टॅब्लो ला अनुमानाचा
कोणताही एक नियम लावल्यावरही
पूर्ततायोग्य टॅब्लो मिळतो.
यावरून प्रमेय सिद्ध होते: कोणताही
बंद टॅब्लो हा~$\Gamma$
मधील गृहीतके एवढीच असलेल्या
टॅब्लो ला नियम लावून मिळतो.
म्हणून $\Gamma$ पूर्ततायोग्य
असेल, तर त्यावरील प्रत्येक
टॅब्लो पूर्ततायोग्य असेल.
पण बंद टॅब्लो पूर्ततायोग्य
नसतो, कारण त्याच्या सर्व शाखा
बंद असतात आणि बंद शाखा
अपूर्ततायोग्य असतात.

आता पूर्ततायोग्य टॅब्लो,
म्हणजे किमान एक पूर्ततायोग्य शाखा
असलेला टॅब्लो, घेऊ.
नियम लावताना एखाद्या शाखेला
चिन्हांकित सूत्रे जोडली जातात
किंवा तिचे दोन शाखांत विभाजन
होते. नियम न लावलेली एखादी
पूर्ततायोग्य शाखा आधीपासून
असेल, तर ती विस्तारित
टॅब्लोमध्येही तशीच
राहते आणि विस्तारित टॅब्लो
पूर्ततायोग्य असतो. म्हणून
पूर्ततायोग्य शाखेलाच नियम
लावल्याचे प्रकरण पाहणे पुरेसे आहे.

त्या शाखेवरील चिन्हांकित सूत्रे
चा संच $\Gamma$ असू द्या आणि
नियम ज्या चिन्हांकित सूत्र
ला लावतो ते
$\sFmla{S}{\varphi}[\sigma] \in \Gamma$
असू द्या. नियमामुळे शाखा फुटत
नसेल, तर नियमाच्या निष्कर्षांसह
$\Gamma$ असलेली विस्तारित शाखा
अजूनही पूर्ततायोग्य आहे हे
दाखवावे लागते. शाखा फुटत
असेल, तर मिळणाऱ्या दोन शाखांपैकी
किमान एक पूर्ततायोग्य आहे हे
दाखवावे लागते.
\begin{enumerate}
\item $\sFmla{\True}{\lnot \psi}[\sigma] \in \Gamma$
  याला $\TRule{\True}{\lnot}$ लावून
  शाखा वाढवू. मग विस्तारित
  शाखेत $\Gamma \cup
  \{\sFmla{\False}{\psi}[\sigma.{*}]\}$
  ही चिन्हांकित सूत्रे असतात.
  $\mSat{M}{\Gamma}[f]$ असे समजा.
  विशेषतः
  $\mSat{M}{\lnot \psi}[f(\sigma)]$.
  म्हणून $Rf(\sigma)w$ असलेल्या
  प्रत्येक $w$ साठी
  $\mSat/{M}{\psi}[w]$; यात
  $f(\sigma.{*})$ सुद्धा येते.
  त्यामुळे $\mModel{M}$ हे
  $\sFmla{\False}{\psi}[\sigma.{*}]$
  $f$ च्या संदर्भात पूर्ण करते.
\item $\sFmla{\False}{\lnot \psi}[\sigma] \in \Gamma$
  याला $\TRule{\False}{\lnot}$
  लावून शाखा वाढवण्याचे प्रकरण:
  सरावासाठी.
\item $\sFmla{\True}{\psi \land \chi}[\sigma] \in \Gamma$
  याला $\TRule{\True}{\land}$
  लावल्यावर त्याच शाखेत
  $\sFmla{\True}{\psi}[\sigma]$
  आणि $\sFmla{\True}{\chi}[\sigma]$
  ही दोन नवी चिन्हांकित सूत्रे
  मिळतात. $\mSat{M}{\Gamma}[f]$
  असे समजा; विशेषतः
  $\mSat{M}{\psi \land
    \chi}[f(\sigma)]$. मग
  $\mSat{M}{\psi}[f(\sigma)]$
  आणि $\mSat{M}{\chi}[f(\sigma)]$.
  म्हणजे $\mModel{M}$ हे
  $\sFmla{\True}{\psi}[\sigma]$
  आणि $\sFmla{\True}{\chi}[\sigma]$
  दोन्ही~$f$ च्या संदर्भात
  पूर्ण करते.
\item $\sFmla{\False}{\psi \lor \chi}[\sigma] \in \Gamma$
  याला $\TRule{\False}{\lor}$
  लावण्याचे प्रकरण: सरावासाठी.
\item $\sFmla{\False}{\psi \lif \chi}[\sigma] \in \Gamma$
  याला $\TRule{\False}{\lif}$
  लावल्यावर शाखेत
  $\sFmla{\True}{\psi}[\sigma.n]$
  आणि $\sFmla{\False}{\chi}[\sigma.n]$
  ही दोन नवी चिन्हांकित सूत्रे
  मिळतात. येथे $\sigma.n$
  हे शाखेसाठी नवे पूर्वचिन्ह आहे: $\sigma.n\notin P(\Gamma)$,
  आणि त्याचा कोणताही विस्तारही $P(\Gamma)$ मध्ये नाही.

  $\Gamma$ पूर्ततायोग्य असल्याने,
  $\mSat{M}{\Gamma}[f]$ असेल
  असे $\mModel{M}$ प्रतिमान
  आणि $P(\Gamma)$ ची~$f$
  अर्थनिर्धारण आहेत. विशेषतः
  $\mSat/{M}{\psi \lif \chi}[f(\sigma)]$.
  पुढील संच पूर्ततायोग्य आहे हे
  दाखवावे लागेल:
  $\Gamma \cup
  \{\sFmla{\True}{\psi}[\sigma.n],
  \sFmla{\False}{\chi}[\sigma.n]\}$.
  त्यासाठी $P(\Gamma)
  \cup \{\sigma.n\}$
  याचे अर्थनिर्धारण पुढीलप्रमाणे
  ठरवू:

  $\mSat/{M}{\psi \lif \chi}[f(\sigma)]$
  असल्याने $Rf(\sigma)w$,
  $\mSat{M}{\psi}[w]$ आणि
  $\mSat/{M}{\chi}[w]$
  असा $w \in W$ असतो.
  $f'$ हे $f$ सारखेच, फक्त
  $f'(\sigma.n) = w$
  असे ठरवू. $f'(\sigma) = f(\sigma)$
  आणि $Rf(\sigma)w$ असल्याने $Rf'(\sigma)f'(\sigma.n)$.
  जुन्या संचातील $\tau$ हा $\sigma.n$ चा आरंभीचा खंड
  असेल, तर तो $\sigma$ चाही आरंभीचा खंड आहे.
  म्हणून जुन्या अर्थनिर्धारणाने $Rf(\tau)f(\sigma)$,
  आणि संक्रमणामुळे $Rf'(\tau)f'(\sigma.n)$.
  नव्या पूर्वचिन्हाचा कोणताही विस्तार जुन्या संचात नाही;
  त्यामुळे त्यापासून जुन्या विस्ताराकडे अतिरिक्त अट येत नाही.
  नव्या पूर्वचिन्हाची स्वतःशी अट परावर्तितेने मिळते.
  म्हणून $f'$ हे $P(\Gamma)\cup\{\sigma.n\}$ चे
  सर्व पूर्वचिन्ह-विस्तारांसाठी अर्थनिर्धारण आहे.
  तसेच $\mSat{M}{\psi}[f'(\sigma.n)]$
  आणि $\mSat/{M}{\chi}[f'(\sigma.n)]$.
  सर्व पूर्वचिन्हे
  $\sigma' \in P(\Gamma)$ साठी
  $f(\sigma') = f'(\sigma')$
  असल्याने
  $\mSat{M}{\Gamma}[f']$.
  म्हणून $\mModel{M}, f'$ हे
  $\Gamma \cup
  \{\sFmla{\True}{\psi}[\sigma.n],
  \sFmla{\False}{\chi}[\sigma.n]\}$
  पूर्ण करते.
\end{enumerate}
आता दोन शाखा निर्माण करणारे
नियम पाहू.
\begin{enumerate}
\item $\sFmla{\False}{\psi \land \chi}[\sigma] \in \Gamma$
  याला $\TRule{\False}{\land}$
  लावल्यावर दोन शाखा मिळतात:
  डावीकडील शाखेत
  $\sFmla{\False}{\psi}[\sigma]$
  आणि उजवीकडील शाखेत
  $\sFmla{\False}{\chi}[\sigma]$.
  $\mSat{M}{\Gamma}[f]$
  असे समजा; विशेषतः
  $\mSat/{M}{\psi \land \chi}[f(\sigma)]$.
  मग $\mSat/{M}{\psi}[f(\sigma)]$
  किंवा $\mSat/{M}{\chi}[f(\sigma)]$.
  पहिल्या प्रकरणात
  $\mModel{M}, f$ हे
  $\sFmla{\False}{\psi}[\sigma]$
  पूर्ण करते; म्हणजे डावी
  शाखा पूर्ततायोग्य आहे.
  दुसऱ्या प्रकरणात
  $\mModel{M}, f$ हे
  $\sFmla{\False}{\chi}[\sigma]$
  पूर्ण करते; म्हणजे उजवी
  शाखा पूर्ततायोग्य आहे.
\item $\sFmla{\True}{\psi \lor \chi}[\sigma] \in \Gamma$
  याला $\TRule{\True}{\lor}$
  लावण्याचे प्रकरण: सरावासाठी.
\item $\sFmla{\True}{\psi \lif \chi}[\sigma] \in \Gamma$
  याला $\TRule{\True}{\lif}$
  लावण्याचे प्रकरण: सरावासाठी.
\end{enumerate}
\end{proof}

\begin{prob}
\ref{int:tab:sou:thm:tableau-soundness}
या प्रमेयाचा पुरावा पूर्ण करा.
\end{prob}

\begin{cor}
\label{int:tab:sou:cor:entailment-soundness}
$\Gamma \Proves \varphi$ असेल,
तर $\Gamma \Entails \varphi$.
\end{cor}

\begin{proof}
  $\Gamma \Proves \varphi$ असेल,
  तर काही
  $\psi_1$, \dots, $\psi_n \in
  \Gamma$ साठी
  $\Delta = \{\sFmla{\False}{\varphi}[1], \sFmla{\True}{\psi_1}[1],
  \dots, \sFmla{\True}{\psi_n}[1]\}$
  याला बंद टॅब्लो असतो.
  $\Gamma \Entails \varphi$
  हे दाखवायचे आहे. तसे नाही
  असे समजा. मग एखाद्या
  $\mModel{M}$ प्रतिमानात
  आणि जगात~$w$,
  $i=1$, \dots,~$n$ साठी
  $\mSat{M}{\psi_i}[w]$,
  पण $\mSat/{M}{\varphi}[w]$.
  $f(1) = w$ असे ठरवू;
  मग $f$ ही
  $P(\Delta)$ ची
  $\mModel{M}$ मधील
  अर्थनिर्धारण आहे आणि
  $\mModel{M}$ हे
  $\Delta$ ला~$f$ च्या
  संदर्भात पूर्ण करते.
  पण \ref{int:tab:sou:thm:tableau-soundness}
  नुसार, बंद टॅब्लो
  असल्याने $\Delta$
  अपूर्ततायोग्य आहे; ही
  विसंगती आहे. म्हणून
  $\Gamma \Entails \varphi$
  हेच मिळते.
\end{proof}

\begin{cor}
\label{int:tab:sou:cor:weak-soundness}
$\Proves \varphi$ असेल, तर
$\varphi$ हे सर्व प्रतिमानांत
सत्य असते.
\end{cor}



\part{प्रतिवास्तविक विधाने}\label{cnt:part}
